\documentclass[10pt]{beamer}

\input{include/preambule.tex}
\makeatletter
\newcounter{lec@number}
\newcounter{lec@globalnumber}
\newcounter{lec@partnumber}

\setcounter{lec@number}{1}
\setcounter{lec@globalnumber}{1}
\setcounter{lec@partnumber}{1}

\edef\@curpart{}

% #1  -  файл лекции
% #2  -  наименование лекци
\NewDocumentCommand{\lec}{smm}{
	\expandafter\edef\csname lec@#2@title\endcsname{#3}
	\expandafter\edef\csname lec@#2@part\endcsname{\@curpart}
	\expandafter\edef\csname lec@#2@partnumber\endcsname{\thelec@partnumber}
	\expandafter\edef\csname lec@#2@number\endcsname{\thelec@number}
	\expandafter\edef\csname lec@#2@globalnumber\endcsname{\thelec@globalnumber}
	\IfBooleanF{#1}
	{
		\stepcounter{lec@number}
		\stepcounter{lec@globalnumber}
	}

}

\NewDocumentCommand{\lecpart}{sm}{
	\edef\@curpart{#2}	
	\setcounter{lec@number}{0}
	\IfBooleanF{#1}
	{
		\stepcounter{lec@partnumber}
	}
	
}
\NewDocumentCommand{\lecturename}{}{%
	\ifcsname lec@\jobname @title\endcsname
		\csname lec@\jobname @title\endcsname
	\else
		---%
	\fi
}
\NewDocumentCommand{\lecturepart}{}{%
	\ifcsname lec@\jobname @part\endcsname
		\csname lec@\jobname @part\endcsname
	\else
		---%
	\fi
}

\NewDocumentCommand{\lecturenumber}{}{%
	\ifcsname lec@\jobname @number\endcsname
		\csname lec@\jobname @number\endcsname
	\else
		---%
	\fi
}

\NewDocumentCommand{\lecturegnumber}{}{%
	\ifcsname lec@\jobname @globalnumber\endcsname
		\csname lec@\jobname @globalnumber\endcsname
	\else
		---%
	\fi
}

\makeatother


\lecpart{Вводный раздел}
\lec{CG_intro}{Введение в курс}
\lec*{CG_history}{История компьютерной графики. Основные вехи}
\lec{CG_vectors}{Математические основы компьютерной графики}
\lec{CG_matrices}{Преобразования точек на плоскости}
\lec{CG_matrices_3d}{Преобразования точек в пространстве}
\lec{CG_rotations}{Вращение пространства}
\lec{CG_pipeline_vertex}{Вершинные операции в графическом конвейере}
\lec{CG_clipping_rasterization}{Отсечение отрезков и растеризация}
\lec{CG_backface_culling}{Удаление невидимых граней}
\lec{CG_lihgtning}{Освещение и закраска объектов}
\lec{CG_textures1}{Текстурирование}
\lec{CG_textures2}{Наложение специальных текстур}




\usepackage{forest}

\expandafter\title{Лекция \lecturenumber \\ \lecturename}
\subtitle{Компьютерная графика}

 \begin{document}
 		 
\makeatletter
\defbeamertemplate*{title page}{my theme}
{
	
	\hfill
	
	\begin{beamercolorbox}[wd=.9\paperwidth,center,]{title}%
		
	\end{beamercolorbox}%	
	
	\vbox to 1em {}
	
	\begin{beamercolorbox}[wd=.9\paperwidth,center,]{title}%
		\usebeamerfont{subtitle}%
		\hfill
		
		\insertsubtitle
		
		\usebeamerfont{title}%
		\inserttitle{} \\[0.5em]
		
		
		
	\end{beamercolorbox}%	
	\hfill\hfill
	
	\begin{beamercolorbox}[wd=.9\paperwidth,center,]{}%
		\usebeamerfont{author}%
		\hfill \\[0.5em]
		\insertauthor{} \\[0.5em]
		\regals
		    
		\vbox to 1em{}
		\usebeamerfont{institute}%
		\insertinstitute {}
		
		\vbox to 1em{}			
		
	\end{beamercolorbox}%	
	\hfill\hfill
	
	
	
	{ \usebeamerfont{institute}\hfill Обновлено: \insertshortdate{} }
	
	\vbox to 5em{}
	
	
}
\defbeamertemplate*{footline}{my theme}{
	\leavevmode%
	\hbox{%
		\begin{beamercolorbox}[wd=.25\paperwidth,ht=3.25ex,dp=0ex,center,sep=1pt]{author in head/foot}%
			\usebeamerfont{author in head/foot}%
			\insertauthor 
			\beamer@ifempty{\insertshortinstitute}{}
		\end{beamercolorbox}%
		\begin{beamercolorbox}[wd=.65\paperwidth,ht=3.25ex,dp=0ex,center,sep=1pt]{title in head/foot}%
			\usebeamerfont{title in head/foot}\insertshortinstitute
		\end{beamercolorbox}%
		\begin{beamercolorbox}[wd=.1\paperwidth,ht=3.25ex,dp=0ex,center, sep=0.5pt]{date in head/foot}%
			\usebeamerfont{date in head/foot}
			\footnotesize \expandafter\toDni{\insertframenumber} / \expandafter\toDni{\inserttotalframenumber}
	\end{beamercolorbox}}%
}

\makeatletter
\beamer@theme@subsectionfalse
\makeatother

 % Head
%\defbeamertemplate*{headline}{miniframes theme}
%{%
%	\begin{beamercolorbox}[colsep=1.5pt]{upper separation line head}
%	\end{beamercolorbox}
%	\begin{beamercolorbox}{section in head/foot}
%		\vskip2pt\insertnavigation{\paperwidth}\vskip2pt
%	\end{beamercolorbox}%
%	\ifbeamer@theme@subsection%
%	\begin{beamercolorbox}[colsep=1.5pt]{middle separation line head}
%	\end{beamercolorbox}
%	\begin{beamercolorbox}[ht=2.5ex,dp=1.125ex,%
%		leftskip=.3cm,rightskip=.3cm plus1fil]{subsection in head/foot}
%		\usebeamerfont{subsection in head/foot}\insertsubsectionhead
%	\end{beamercolorbox}%
%	\fi%
%	\begin{beamercolorbox}[colsep=1.5pt]{lower separation line head}
%	\end{beamercolorbox}
%}

{
	
}


\makeatother






%float page top aligment
\makeatletter
\setlength{\@fptop}{0pt}
\setlength{\@fpbot}{0pt plus 1fil}
\makeatother
 		 


	\newcommand \abs[1] {\left| #1 \right|}
	
	\everymath{\displaystyle}  
  
    \QRFRAME	
    
	\frame{\maketitle}

	\begin{frame}{План лекции}
		\tableofcontents
	\end{frame}
	
	\begin{frame}{Виды алгоритмов удаления невидимых граней и поверхностей}
		
		\begin{forest}			
			[Алгоритмы удаления\\невидимых граней и поверхностей, for tree={draw, align=center, minimum height=3em}
			[Работающие в\\пространстве экрана]
			[Работающие в\\пространстве объектов]
			[Комбинированные варианты\\~]
			]			
		\end{forest}
		
	\end{frame}
		
	\section{Метод плавающего горизонта}
	
	\begin{frame}{Метод плавающего горизонта}
		\centering%
		\includegraphics[width=0.8\textwidth]{CreateSurfacePlotExample_01.png}
	\end{frame}
	
	\begin{frame}{Метод плавающего горизонта}
		\TC{0.5}
		{
			\begin{enumerate}
				\item \includegraphics[width=0.8\textwidth]{mpg_1.png}
				\item \includegraphics[width=0.9\textwidth]{mpg_2.png}
			\end{enumerate}
		}
		{
			\begin{enumerate}
				\setcounter{enumi}{2}
				\item \includegraphics[width=0.9\textwidth]{mpg_3.png}
				\item \includegraphics[width=0.9\textwidth]{mpg_4.png}
			\end{enumerate}
			
		}
		
	\end{frame}
	
	\begin{frame}{Примеры}
		
		\TC{0.5}
		{
			\centering%
			Пример хорошей отрисовки
			
			\fbox{\animategraphics[autoplay,loop, nomouse, width=0.85\textwidth]{5}{Images/\jobname/nikulin_good_draw/nikulin_good_draw-}{0}{69}}
		}
		{
			\centering%
			Пример плохой отрисовки
			\fbox{\animategraphics[autoplay,loop, nomouse, width=0.85\textwidth]{5}{Images/\jobname/nikulin_bad_draw/nikulin_bad_draw-}{0}{44}}
		}
		
		~
		
		\HREF{https://grafika.me/node/1093}{Демонстрация}
		
		
		
	\end{frame}
	

	
		
	\section{Отсечение нелицевых граней}
	
	
	\begin{frame}{Отсечение нелицевых граней}
		
	
			
			%\draw[very thin] (-5,-5) grid[step=0.5] (5,5);
%			\fill [red] circle (5pt);
%			\foreach \t in {-5,...,5}
%				{
%					\draw (\t,0) node {\t};
%					\draw (0,\t) node {\t};
%				}

			\TC{0.4}
			{
			\centering%
			\begin{tikzpicture}
						\coordinate (A) at (0,0);
			\coordinate (B) at (2,-1);
			\coordinate (C) at (3,0.2);
			\coordinate (D) at (0.9, 2);
			
			\draw (A) -- (B) -- (C) -- cycle;
			\draw (A) -- (D);
			\draw (B) -- (D);
			\draw (C) -- (D);
			
			\coordinate (G1) at (1,0.5);
			
			\draw [rotate=80, shift=(G1)]  circle[x radius = 0.3, y radius = 0.2];
			\draw[-Latex] (G1) -- +(170:1);
			
			\coordinate (G2) at (2,0.5);
			
			\draw [rotate=130, shift=(G2)]  circle[x radius = 0.3, y radius = 0.2];
			\draw[-Latex] (G2) -- +(20:1);
			
			\coordinate (G3) at (1.5,-0.3);
			
			\draw [rotate=0, shift=(G3),dashed]  circle[x radius = 0.3, y radius = 0.2];
			\draw[-Latex] (G3) -- +(-90:1);
			
			
			\coordinate (G4) at (1.2,1.1);
			
			\draw [rotate=0, shift=(G4),dashed]  circle[x radius = 0.25, y radius = 0.3];
			\draw[-Latex] (G4) -- +(120:0.7);
			
			\end{tikzpicture}
		}
		{
			
				\textbf{Для параллельной проекции}
				
					Для каждой грани $f$ необходимо проверить знак $\vv n \cdot \vv v$. 
					
					
					$\vv n \perp f$, $\vv v$ - вектор на наблюдателя ($\perp$ к плоскости проецирования)
				
				~
				
				
				\textbf{Для центральной проекции}
				
				Для каждой грани $f$ необходимо проверить знак $\VV n \cdot \VV{(c-p)}$. 
				
				$p \in f$, $c$ --- центр проекции.
	
		}
				

		
	\end{frame}
	
	\section{Алгоритм Робертса}
	
	\begin{frame}{Алгоритм Робертса}
		\centering%
		\includegraphics[width=0.9\textwidth]{roberts.png}
	\end{frame}
	
		\begin{frame}{Оптимизация алгоритма Робертса}
		\centering%
		\includegraphics[width=0.9\textwidth]{roberts_opt2.png}
	\end{frame}
	
		\begin{frame}{Оптимизация алгоритма Робертса}
		\centering%
		\includegraphics[width=0.9\textwidth]{roberts_opt.png}
	\end{frame}
	
	
	\section{Алгоритм Уорнока}
	
	\begin{frame}{Алгоритм Уорнока}
		\TC{0.45}
		{
			\centering%
			\includegraphics[width=1\textwidth]{varnak.png}	
		}
		{
			Случаи, когда можно легко определить видимость:
			\begin{itemize}				
			
			\item окне нет ни одной грани, в этом случае ничего рисовать не нужно;
			\item в окне ровно одна грань, в этом случае достаточно отобразить эту
			грань;
			\item размеры окна равны размерам одного пиксела, в этом случае из
			граней выбирается ближайшая;
			\item ближайшая грань охватывает всё окно, в этом случае она заслоняет все остальные грани.
			
			\end{itemize}
		}
		
		Алгоритм можно разделить на следующие элементарные задачи:\\
		- определение, в какие части окна попадает грань;\\
		- определение, что грань охватывает окно;\\
		- определение, что грань является ближайшей.
	\end{frame}
	
	\begin{frame}{Разделение окна вдоль общего ребра двух граней}
		\centering%
		\includegraphics[width=0.7\textwidth]{varnak2.png}	
	\end{frame}
	
	\section{Алгоритм Z-буфера}
	
	\begin{frame}{Алгоритм Z-буфера}
		\TC{0.65}
		{
			\centering\includegraphics[width=0.95\textwidth]{zbuffermatrix.png}
		}
		{
			\centering\includegraphics[width=0.95\textwidth]{Z-buffer_example.jpg}
		}
	\end{frame}
	
	\section{Алгоритм художника}
	
	\begin{frame}{Алгоритм художника}
		\TC{0.7}		
		{\centering\includegraphics[width=\textwidth]{artist.png}}		
		{\centering\includegraphics[width=\textwidth]{artist2.png}}
		
	\end{frame}
	
	
		\begin{frame}{Алгоритм художника}
	\centering\includegraphics[width=0.9\textwidth]{artist3.png}		
	
		
	\end{frame}
	\note{
		\footnotesize
	Одним из методов сортировки граней является метод бинарного раз-
	деления пространства. В качестве исходных данных выступает набор гра-
	ней в пространстве и точка, в которой находится наблюдатель. Если разде-
	лить пространство плоскостью одной из граней, то грани, лежащие в полу-
	пространстве, в котором находится наблюдатель, будут ближе к нему, чем
	те, что лежат в другом полупространстве. Поэтому порядок рисования гра-
	ней будет следующим:
	
	- грани, лежащие в полупространстве, где нет наблюдателя;
	
	- грани, лежащие на разделяющей плоскости;
	
	- грани, лежащие в полупространстве, где находится наблюдатель.
	
	Аналогично сортируются и грани в полупространствах, в каждом выбира-
	ется грань, её плоскостью полупространство делится на две части и так до
	тех пор, пока в каждой части не останется ни одной грани.
	
	В качестве примера рассмотрим случай на плоскости (см. рис. 15),
	вместо граней – отрезки (a, b, c, d, e, f, g, h), вместо плоскостей – прямые
	(можно представить, что все грани вертикальные, а на рисунке показан их
	вид сверху).
	
	Наблюдатель находится в точке V. Плоскость первой грани (a) делит
	грани на три группы: (b, c, d) – лежат в полупространстве наблюдателя, (e,
	g, h, f) – в другом полупространстве, грань (a) – на разделяющей плоско-
	сти. Порядок рисования (e, g, h, f), (a), (b, c, d). Грани в каждом из полу-
	пространств нужно отсортировать межу собой. В одном выбирается плос-
	кость грани (e), она делит грани на группы: (f1) – лежит в полупространст-
	ве наблюдателя, (f2, g) – в другом полупространстве, грани (e, h) – на раз-
	деляющей плоскости. Грань (f) попала в оба полупространства и была раз-
	делена на две части (f1 и f2), грань (h) попала на разделяющую плоскость и должна рисоваться вместе с гранью (e), их взаимный порядок рисования
	может быть любым, они не могут заслонить друг друга, так как лежат в
	одной плоскости. Аналогично, в части пространства с гранями (f2, g) выби-
	рается плоскость g, в части пространства с гранями (b, c, d) – грань (c). По-
	сле всех разделений порядок рисования граней будет: f2, g, e, h, f1, a, d, c, b.
	
	
	}
	
	\section{Алгоритм трассировки лучей}
	
	\begin{frame}{Алгоритм трассировки лучей}
		\centering\includegraphics[width=0.9\textwidth]{raytrace.png}		
		
		
	\end{frame}
	
	\section{Алгоритм построчного сканирования}
	
	\begin{frame}{Алгоритм построчного сканирования}
		\centering\includegraphics[width=0.9\textwidth]{string.png}			
		
	\end{frame}
	
	\section{Алгоритм S-буфера}
	
	\begin{frame}{Алгоритм S-буфера}
		\centering\includegraphics[width=0.9\textwidth]{sbuffer1.png}		
	\end{frame}
	
	\begin{frame}{Разбиение отрезков}
		\centering\includegraphics[width=0.9\textwidth]{sbuffer2.png}		
	\end{frame}
	


	
 
 
\begin{comment}
	
\end{comment}
 

\end{document}